% ============================================================================
\section{Méthodes d'ordonnancement}
\label{sec:algos}
% ============================================================================

Tous les algorithmes utilisent le même modèle (section~\ref{sec:formulation}) et
produisent un \code{allocation.csv} directement exécutable par l'orchestrateur.

\subsection{Algorithmes de référence}

Ce sont les heuristiques classiques de \emph{list scheduling} sur machines
hétérogènes ($n$ tâches, $m$ VM).

\begin{table}[H]
\centering\small
\begin{tabularx}{\textwidth}{>{\raggedright\arraybackslash}p{2.1cm}X>{\raggedright\arraybackslash}p{1.7cm}}
\toprule
\textbf{Algorithme} & \textbf{Principe} & \textbf{Coût} \\
\midrule
Round Robin & tâche $i$ $\rightarrow$ VM $i \bmod m$ ; ignore durées, prix et charge. \textbf{Référence} ($S = 1$). & $O(n)$ \\
MET & chaque tâche va sur la VM où \emph{sa} durée est minimale, sans regarder la file. & $O(nm)$ \\
MCT & dans l'ordre du lot, chaque tâche va sur la VM où elle \emph{finirait} le plus tôt, compte tenu des tâches déjà placées. & $O(nm)$ \\
Min-Min & à chaque tour, on place la tâche dont la meilleure heure de fin est la plus \emph{petite}. & $O(n^2 m)$ \\
Max-Min & à chaque tour, on place la tâche dont la meilleure heure de fin est la plus \emph{grande} (grosses tâches d'abord). & $O(n^2 m)$ \\
\bottomrule
\end{tabularx}
\caption{Heuristiques de référence (\code{scheduler/algorithms/}).}
\end{table}

Ces heuristiques optimisent le \textbf{temps} : aucune ne regarde le prix des VM,
et elles utilisent toujours toutes les machines disponibles. Les simulations de
la section~\ref{sec:simu} montrent que c'est précisément ce qui leur manque.

\subsection{Notre algorithme : \sel}

\paragraph{Idée.} Les mesures ont montré que (1) les VM ont la même vitesse par
tâche, (2) une VM allumée coûte au moins 60~s, et (3) Max-Min répartit déjà bien
les tâches. Le levier principal n'est donc pas la répartition fine, mais
\textbf{le choix des VM à allumer}. De plus, deux VM du même type (vm2 et vm3)
sont interchangeables : on ne se demande pas « quelles VM allumer », mais
\textbf{« combien de VM de chaque type allumer »}.

\paragraph{Méthode, en deux niveaux.}
\begin{enumerate}
  \item \textbf{Choix du nombre de VM par type — exact.} On essaie toutes les
        combinaisons : de 0 à $n_t$ VM pour chaque type $t$, sauf « aucune VM ».
        Pour notre parc : e2-standard-2 (0–1) $\times$ e2-medium (0–2) $\times$
        e2-highcpu-4 (0–1), soit $2 \times 3 \times 2 - 1 = 11$ combinaisons.
  \item \textbf{Répartition — heuristique.} Sur chaque combinaison, Max-Min
        (le meilleur algorithme classique dans nos simulations).
\end{enumerate}
On garde la combinaison de plus petit score $S$ parmi celles qui respectent le
délai et la mémoire.

\begin{algorithm}[H]
\caption{\sel}
\begin{algorithmic}[1]
\Require tâches $T$, VM $M$ regroupées par type, modèle $p_{ij}$
\Ensure une allocation ordonnée
\State $\mathit{ref} \gets$ évaluation de Round Robin sur toutes les VM \Comment{donne $\Drr$, $\Crr$}
\State $\mathit{meilleure} \gets$ aucune
\For{chaque combinaison $(c_1,\dots,c_K)$ avec $0 \le c_t \le n_t$, non toutes nulles}
  \State $V \gets$ les $c_t$ premières VM de chaque type $t$
  \State $A \gets$ Max-Min restreint aux VM de $V$
  \State calculer $D(A)$, $C(A)$ par simulation, puis $S(A)$
  \If{$D(A) \le 1{,}25\,\Drr$ \textbf{et} RAM respectée \textbf{et} $S(A) < S(\mathit{meilleure})$}
    \State $\mathit{meilleure} \gets A$
  \EndIf
\EndFor
\State \Return $\mathit{meilleure}$ (ou Round Robin si aucune combinaison n'est admissible)
\end{algorithmic}
\end{algorithm}

\paragraph{Justification et propriétés.}
\begin{itemize}
  \item \textbf{Décomposition} : choisir les ressources, puis ordonnancer dessus.
        Chaque partie s'explique simplement ; la première est exacte.
  \item \textbf{Jamais pire que Max-Min} (dès que Max-Min est admissible) : la
        combinaison « toutes les VM » fait partie des essais. Si aucune combinaison
        n'est admissible, on renvoie Round Robin, admissible par définition.
        Ces propriétés sont vérifiées par des tests automatiques.
  \item \textbf{Général} : aucun nom de VM, prix ou durée n'est écrit dans le code ;
        tout vient de \code{machines.csv} et \code{vm\_performance.csv}. Le même code
        éteint 2, 1 ou 0 VM selon le lot (section~\ref{sec:simu}). L'objectif reste
        paramétrable (\code{-wcost}, \code{-deadline}).
  \item \textbf{Optimise directement $S$}, contrairement à notre première version,
        dont la recherche locale minimisait le makespan sous un budget de coût.
  \item \textbf{Déterministe et explicable} : \code{./bin/scheduler -explain} affiche
        les 11 combinaisons et leur score (tableau~\ref{tab:explain500}).
\end{itemize}

\paragraph{Complexité et passage à l'échelle.} Le nombre de combinaisons est
$\prod_t (n_t + 1) - 1$ : il dépend du nombre de \emph{types}, pas du nombre de
VM. Avec 3 types et 10 VM de chacun, on obtient 1\,330 combinaisons, contre
environ $10^9$ si l'on traitait chaque VM séparément. Chaque combinaison coûte
un Max-Min en $O(n^2 m)$. Temps mesuré : \textbf{114~ms pour 2000 tâches}. Au-delà
de quelques types, il faudrait remplacer l'énumération par un glouton (ajouter
les types au meilleur débit par dollar jusqu'à tenir le délai).

\paragraph{Intégration.} En mode \code{-algo \sel}, le scheduler affiche aussi
les VM à démarrer (\code{VMS="..." ./scripts/manage\_vms.sh start}) : les VM non
choisies restent éteintes.
